\chapter{DOGMA state construction and execution}
\label{ch:evo31}
\section{Make the native state executable}
Chapter 30 supplied a serialized request lifecycle: admission, ownership,
capacity reservation, completion fencing and cleanup. It did not specify how
a native model consumes a token. This chapter fills that narrow gap with a
fully auditable recurrence, then attaches its state changes to the shell.

Our implementation is an \Term{input-gated diagonal recurrence}, not a complete
DOGMA model. DOGMA denotes the non-Transformer research direction; HermonDNA
denotes the Transformer branch; Evolutor is the broader orchestration and research
system. The toy here establishes arithmetic and interface properties, not
training quality, genomic capability, architectural novelty or production
performance. Its gate depends on current input, not previous hidden state.
This restriction is precisely why an independent affine closed form is possible.

The reader will construct a content-bound state, ingest a masked prompt without
emitting output, calculate logits by hand, verify derivatives independently,
and stage a decode step that cannot advance native state after cancellation.
All tables and literal code listings come from
\texttt{drafts/ch31/reference.py}.

\section{Declare the artifact, not only its name}
The artifact contains an ordered vocabulary of size $V$, embedding
$E\in\mathbb R^{V\times D}$, retention $a\in[0,1)^D$, write parameter
$w\in\mathbb R^D$ and output head $H\in\mathbb R^{D\times V}$.
The reference requires finite CPU float64 tensors and an explicit schema version.
Reordering vocabulary strings while leaving the matrix unchanged changes the
meaning of every token ID.

A \Term{content-bound artifact} is identified by canonical serialization of
schema, dtype, device, vocabulary and all parameter values. The SHA-256 digest
is recomputed at execution boundaries in this teaching oracle.
It detects changes, including in-place tensor edits, rather than trusting a
human-readable model label. It does not authenticate the publisher, prove that
a checkpoint was trained, or certify that a hidden state was reached by the
advertised prefix.

Production systems should validate immutable loaded artifacts once and use
protected handles. Hashing all parameters on every token is intentionally
expensive here and is not a proposed high-throughput serving design.
\Plate{01-constructor}{Artifact construction binds data representation and token semantics. A frozen Python record does not freeze the storage of its tensors.}{fig:evo31-constructor}
\Listing{artifact_id}

\section{Construct and validate the state}
The state header is
\[
 S=(\text{tenant},\text{request},\text{artifact digest},p,h),
 \qquad p\in\mathbb N_0,\quad h\in\mathbb R^D.
\]
Construction initializes $p=0$ and $h=\mathbf0$ without consuming a prompt.
Tenant and request form an exact ownership pair. A valid header, correctly
shaped finite hidden vector and matching artifact digest are checked on the
actual native-step path, including an empty ingestion call.

This is stronger than an optional assertion that callers may forget to invoke.
It is weaker than authentication: trusted Python code can construct or alter
objects. State reachability, signed checkpoint provenance and authorization are
separate obligations. The reference also has no GPU/device-transfer policy,
mixed precision, batch allocator or persisted checkpoint format.

\section{Derive one native step}
For a token ID $x_t\in\{0,\ldots,V-1\}$,
\begin{align}
 e_t&=E[x_t]\in\mathbb R^D,\\
 q_t&=\sigma(w\odot e_t)\in(0,1)^D,\\
 b_t&=q_t\odot e_t,\\
 h_t&=a\odot h_{t-1}+b_t,\\
 z_t&=h_tH\in\mathbb R^V.
 \label{eq:evo31-step}
\end{align}
$\odot$ denotes elementwise multiplication. The gate weights the current
embedding's write, not the historical state. Signed embeddings still produce
signed writes; a sigmoid gate does not turn $h$ into probabilities.

The state position counts consumed live tokens, not predicted tokens or
delivered bytes. Each accepted native step produces a new tensor and increments
$p$ once. No input tensor is updated in place. A nonfinite hidden vector or
logit vector is rejected even when all parameters were initially finite.
\Plate{03-step}{Shape-annotated native dataflow. Retention transports history; the input-only gate contributes a new write. There is no attention matrix or KV cache.}{fig:evo31-step}
\Listing{native_step}

\subsection{Bound state magnitude without promising useful memory}
If $|b_{t,j}|\le M_j$, then repeated substitution yields
\[
 |h_{t,j}|\le a_j^t|h_{0,j}|+
 M_j\sum_{k=0}^{t-1}a_j^k
 =a_j^t|h_{0,j}|+M_j\frac{1-a_j^t}{1-a_j}.
\]
For this finite vocabulary, $M_j=\max_v|E_{v,j}|$ is sufficient because
$0<q_{t,j}<1$. The bound can become large when $a_j$ approaches one.
It bounds magnitude, not the semantic quality or recoverability of remembered
information. Fixed $D$ also does not imply constant execution cost: the dense
output projection costs $O(DV)$ per live token, and this oracle's digest checks
read the complete parameter artifact.

\section{Ingest chunks and masks}
The ingestion function validates all live token IDs before entering the loop.
A malformed tail therefore cannot mutate the caller's input state. It also
requires a Boolean mask with exactly one flag per token row.
A false row performs no embedding lookup, holds the hidden state and position,
and emits an all-zero row vector by convention. That vector is a placeholder,
not a probability distribution or a next-token prediction.

Masking is not request reset. A new request requires fresh construction and a
new owner; silently resetting at padding would corrupt continuation semantics.
An empty chunk retains the state but still checks owner and artifact.
\Plate{02-ingest}{Masked chunk example. Padding neither consumes a token nor advances position. The full chunk is preflighted before live-row execution.}{fig:evo31-ingest}
\Listing{ingest}

\section{Verify exact arithmetic before optimizing}
Choose vocabulary strings $(\texttt{0},\texttt{1},\texttt{2})$ and
\[
 E=\begin{pmatrix}0&0\\1&0\\0&1\end{pmatrix},\quad
 a=(1/2,1/4),\quad w=(0,0),\quad
 H=\begin{pmatrix}0&1&-1\\0&-1&1\end{pmatrix}.
\]
Every gate equals $1/2$. Starting from zero, the prefix $(1,2,1)$ gives:
\begin{center}\input{\Generated/trace.tex}\end{center}
For example, the second state is
$(\tfrac12\cdot\tfrac12,\tfrac14\cdot0+\tfrac12)
=(\tfrac14,\tfrac12)$.
The third is $(\tfrac58,\tfrac18)$, and its logits are
$(0,\tfrac12,-\tfrac12)$.
These rational values provide an independent oracle stronger than comparing
two calls to the same implementation.
\Plate{05-output}{An exact untrained fixture. Consuming a prefix and predicting a following token are distinct operations. Values are computed arithmetic, not a genomic benchmark.}{fig:evo31-output}

\section{A closed form and ordered chunk composition}
Since $b_t$ depends only on the token,
\begin{equation}
 h_{T,j}=a_j^T h_{0,j}+
       \sum_{i=1}^{T}a_j^{T-i}\sigma(w_jE_{x_i,j})E_{x_i,j}.
 \label{eq:evo31-closed}
\end{equation}
The scalar implementation expands this sum using Python arithmetic.
It never calls \texttt{native\_step}; therefore it can detect indexing,
retention and head-projection mistakes in the tensor path.
\Listing{scalar_closed_form}

A chunk represents $h\mapsto A\odot h+B$. Applying left chunk then right chunk
gives
\[
 (A_R,B_R)\circ(A_L,B_L)=
 (A_R\odot A_L,\ A_R\odot B_L+B_R).
\]
Composition is associative because these are affine functions. It is generally
not commutative: changing token order changes the accumulated writes.
This recovers the restricted affine algebra of Chapter 26, not proof that a
full state-dependent DOGMA implementation supports the same scan.
\Plate{04-chunks}{Ordered affine chunk composition and gradient transport. Forward parity, derivative parity and storage independence are separate properties.}{fig:evo31-chunks}

\section{Check derivatives independently}
For a scalar final-state component,
\[
 \frac{\partial h_{T,j}}{\partial w_j}=
 \sum_{i=1}^{T}a_j^{T-i}
 \sigma(w_jE_{x_i,j})(1-\sigma(w_jE_{x_i,j}))E_{x_i,j}^{\,2}.
\]
At $w=0$ for prefix $(1,2,1)$, this gives
$\partial(h_{T,1}+h_{T,2})/\partial w=
(\tfrac14(1+\tfrac14),\tfrac14\cdot\tfrac14)
=(.3125,.0625)$.
The tests compare autograd with both this derivation and central finite
differences at multiple write values.

For the initial state,
$\partial h_{T,j}/\partial h_{0,j}=a_j^T$.
A small $a_j$ can erase gradients even though forward values remain bounded.
Near-one retention trades slower forgetting for a larger magnitude envelope.
Neither observation says that the model has learned useful long-range
dependencies.

\subsection{Copying, detaching and aliasing are not synonyms}
PyTorch distinguishes a differentiable copy from detachment
\cite{evo31-clone,evo31-detach}. A clone preserves the graph connection.
Detachment removes that connection but shares underlying storage.
Consequently, replacing a chunk boundary with \texttt{detach()} can preserve
forward outputs while breaking training parity, and may still permit aliasing.
A detached, storage-independent inference snapshot requires an explicit copy;
a training continuation must retain the appropriate graph instead.

The chunk tests split the prefix at every boundary and compare the derivative
against the analytical result. This does not verify gradients of a larger
learned architecture or its optimizer.

\section{Define next-token semantics}
$z_t$ predicts a token following the consumed prefix through $t$.
Greedy decoding chooses $\hat x_{t+1}=\arg\max z_t$ and then consumes that
selected token to prepare the next prediction. For the fixture after
$(1,2,1)$, greedy ID $1$ is new output; it is not an instruction to re-emit the
last prompt token because both happen to be equal.

The adapter's state includes every generated token after its successful commit.
Thus $p=p_{\rm prompt}+r$, where $r$ is the shell's committed decode revision.
The fixture uses one-character token strings; it has no EOS convention, trained
tokenizer, stochastic sampler or multi-token text detokenizer.
A session requires an already consumed prompt or an explicitly supplied BOS
token. It does not invent a BOS ID or silently treat zero-state logits as a
prompt-conditioned prediction.

\section{Stage the native result before shell publication}
Chapter 30's \texttt{Runtime.launch} reserves queue capacity and returns a ticket
bound to a lease generation and revision. Our \Term{staged decode} computes a
candidate result privately. Native state remains unchanged until
\texttt{Runtime.commit} accepts the corresponding output.

The adapter checks owner, artifact, current state stamp, candidate stamp,
position and head consistency. The stamps detect accidental tensor mutation
while work is pending. Parameter changes are rejected by the digest.
Duplicate or foreign completion objects are rejected before publication.
\Plate{06-shell}{Serialized in-process commit ordering. Late or cancelled work settles its ticket but never installs the candidate hidden state.}{fig:evo31-shell}
\Listing{Session.finish}

Three distinct outcomes matter:
\begin{enumerate}
\item Capacity unavailable: no ticket, no native computation, no state advance.
\item Valid accepted completion: commit one output item, then install its
already validated candidate state.
\item Cancelled, deadline-expired or invalid completion: settle/discard work
without installing the candidate.
\end{enumerate}
For a live prompt state at position three and zero decode revision, a successful
step produces position four and revision one. A cancelled step leaves position
three and revision zero. The tests check both facts, not just the terminal reason.

This is a serialized, trusted, in-process adapter. Publication and Python
assignment are not a durable distributed transaction: a process crash between
them requires recovery machinery that is absent here. The adapter also does not
schedule prefill, run GPU kernels, authenticate callers, protect Python memory
from hostile code, or implement a network transport.

\section{What the tests establish}
All length-three fixture sequences are checked against independent scalar
arithmetic. Additional tests exercise mask semantics, invalid tails, vocabulary
reordering, parameter mutation, wrong owners, malformed states, derivative
continuity, finite differences, cancellation, deadlines, backpressure and
duplicate completion. Suffix changes do not alter earlier prefix outputs.

These are bounded deterministic checks of a declared teaching implementation.
They do not measure throughput, prove security against a malicious caller,
demonstrate biological validity or establish equivalence to a trained DOGMA
model. Trace metadata should not log token IDs or hidden values in an untrusted
production environment; the printed fixture is deliberately public toy data.

\section{Worked exercises}
\begin{enumerate}
\item \textbf{Shapes.} For $V=3,D=2$, give embedding and head shapes.
\emph{Answer:} $[3,2]$ and $[2,3]$.
\item \textbf{Gate.} What is $q$ when $w=0$? \emph{Answer:} $1/2$ in each
coordinate, independently of the embedding.
\item \textbf{First token.} Consume ID 1 from zero. \emph{Answer:}
$h=(.5,0)$ and $z=(0,.5,-.5)$.
\item \textbf{Second token.} Then consume ID 2. \emph{Answer:}
$h=(.25,.5)$ and $z=(0,-.25,.25)$.
\item \textbf{Order.} Reverse the prefix $(1,2)$.
\emph{Answer:} $(2,1)$ gives $h=(.5,.125)$, demonstrating noncommutativity.
\item \textbf{Magnitude.} With $M=.5,a=.5,h_0=0$, give the asymptotic bound.
\emph{Answer:} $M/(1-a)=1$.
\item \textbf{Semantics.} Are signed hidden values probabilities?
\emph{Answer:} no; neither the embedding nor output projection is normalized.
\item \textbf{Masked row.} What does a false mask do?
\emph{Answer:} no lookup, unchanged hidden state and position, zero placeholder
logits, not a prediction.
\item \textbf{Invalid tail.} Does a valid first row authorize consuming an
invalid second row? \emph{Answer:} no; full live-ID preflight rejects the chunk.
\item \textbf{Owner.} Can another tenant use the same request string?
\emph{Answer:} not with this state; the ownership pair must match exactly.
\item \textbf{Empty chunk.} Can it bypass artifact validation?
\emph{Answer:} no; validation occurs before the loop.
\item \textbf{Vocabulary.} Does swapping strings preserve artifact identity?
\emph{Answer:} no; ordered strings are part of canonical content.
\item \textbf{Mutation.} Does a frozen dataclass freeze its tensors?
\emph{Answer:} no; storage remains mutable, so the oracle rechecks its digest.
\item \textbf{Closed form.} Find the coefficient transporting write $b_1$
through a length-three prefix. \emph{Answer:} $a^2$.
\item \textbf{Derivative.} Find the fixture's write gradient.
\emph{Answer:} $(.3125,.0625)$ for the sum of final hidden components.
\item \textbf{Graph.} Can detachment preserve forward values but break training?
\emph{Answer:} yes; it removes the gradient edge.
\item \textbf{Storage.} Does detachment produce independent storage?
\emph{Answer:} no; use an explicit copy when independence is required.
\item \textbf{Position.} A prompt consumes three tokens and one decode commits.
\emph{Answer:} native position four, shell decode revision one.
\item \textbf{Cancellation.} A staged step completes after cancellation.
\emph{Answer:} discard the candidate; native position and committed revision
remain unchanged.
\item \textbf{Durability.} Is local commit-plus-assignment a crash-safe
transaction? \emph{Answer:} no; a persistent recovery protocol is a separate
implementation obligation.
\end{enumerate}

\section{Handoff to state pools and memory ownership}
Chapter 32 makes logical state reusable in bounded physical storage.
Its generational handles, scrubbing and clone semantics must preserve the
owner, artifact, position and hidden-state contract established here. An
allocator slot is not a request identity. The gradient distinction above also
constrains whether a pool copy is an inference snapshot or a training continuation.
